\chapter*{空间向量}

\section*{空间向量及其运算}

\begin{enumerate}

\item 下列说法中正确的是\blkda{(1)(4)(5)}.
    \begin{enumerate}
        \item 空间任意两个共起点的向量是共面的；
        \item 两个相等向量就是相等长度的两条有向线段表示的向量；
        \item 两个相等向量就是起点相同终点也相同的向量；
        \item 空间向量满足加法交换律和结合律；
        \item 首尾相接的若干向量之和等于由起始向量的起点指向末尾向量的终点的向量。
        
    \end{enumerate}

    \item 下列命题正确的是\dotfill{(\kh{C})}

    \xx{若 $\left| \vv{a}\right|  = \left| \vv{b}\right|$ ，则 $\vv{a} = \vv{b}$ }
    {若 $\left| \vv{a}\right|  > \left| \vv{b}\right|$ 则 $\vv{a} > \vv{b}$}
    {若 $\vv{a} = \vv{b}$ ,则 $\left| \vv{a}\right|  = \left| \vv{b}\right|$ }
    {若 $\vv{a} \neq  \vv{b}$ ，则 $\vv{a}$ 与 $\vv{b}$ 的方向不同}


    \item 正方体 ${ABCD} - {A}_{1}{B}_{1}{C}_{1}{D}_{1}$ 中，下列各式中结果为 $\vv{B{D}_{1}}$ 的是\blkda{(1)(2) }
    \begin{enumerate}
        \item $\left( {\vv{{A}_{1}{D}_{1}} - \vv{{A}_{1}A}}\right)  - \vv{AB}$ ；
        \item $\left( {\vv{BC} + \vv{B{B}_{1}}}\right)  - \vv{{D}_{1}{C}_{1}}$ ；
        \item $\left( {\vv{AD} - \vv{AB}}\right)  - 2\vv{D{D}_{1}}$ ；
        \item  $\left( {\vv{{B}_{1}{D}_{1}} + \vv{{A}_{1}A}}\right)  + \vv{D{D}_{1}}$ 。
    \end{enumerate}


    \item 若正四面体 ${ABCD}$ 的棱长为 1， $M$ 是 ${AB}$ 中点，则 $\vv{MC} \cdot  \vv{MD} =$ \blkda{$\dfrac{1}{4}$}。

    \item 已知线段 ${AB},{BD}$ 在平面 $A$ 上， ${BD} \bot  {AB}$ ，垂足为点 $B$ ，线段 ${AC} \bot$ 平面 $C$ ，垂足为 $A$ ，若 ${AB} = a,{BD} = b,{AC} = c$ 。 $C,D$ 之间的距离是\blkda[2]{$\sqrt{a^2+b^2+c^2}$}。

    \item 平行六面体 ${ABCD} - {A}_{1}{B}_{1}{C}_{1}{D}_{1}$ 的体对角线 ${B}_{1}D$ 与 ${A}_{1}B{C}_{1}$ 的交点是 $\bigtriangleup  {A}_{1}B{C}_{1}$ 的\blkda{重}心。

    \item 已知 $\left| \vv{a}\right|  = 2,\left| \vv{b}\right|  = 3, < \vv{a},\vv{b} >  = {60}^{ \circ  }$ ，则 $\left| {2\vv{a} - 3\vv{b}}\right|$ 等于\blkda{$\sqrt{61}$}


    \item 已知在平行六面体 ${ABCD} - {A}_{1}{B}_{1}{C}_{1}{D}_{1}$ 中，同一顶点为端点的三条棱长都等于 1，且彼此的夹角都是 ${60}^{ \circ  }$ ，则此平行六面体的对角线 $A{C}_{1}$ 的长为\blkda{$\sqrt{6}$}


    \item 已知 $a,b$ 是两异面直线， $A,B \in  a,C,D \in  b,{AC} \bot  b,{BD} \bot  b$ 且 ${AB} = 2,{CD} = 1$ ，则直线 $a,b$ 所成的角为\blkda{$\dfrac{\pi}{3}$}


    \item 已知 $A\text{ 、 }B\text{ 、 }P$ 三点共线， $O$ 为空间任意一点， $\vv{OP} = \alpha \vv{OA} + \beta \vv{OB}$ ，则 $\alpha  + \beta  =$ \blkda{1}.

    \item 已知空间四边形 ${ABCD}$ 中， $\vv{AB} = \vv{a} - 2\vv{c},\vv{CD} = 5\vv{a} + 6\vv{b} - 8\vv{c}$ ，对角线 ${AC}$ 、 ${BD}$ 的中点分别为 $E,F$ ,则 $\vv{EF} =$ \blkda[4]{$3\vv{a}+3\vv{b}-5\vv{c}$}.

    \item 已知 $A,B,C$ 三点共线，则对空间任一点 $O$ ，存在三个不为 0 的实数 $\lambda ,m,n$ ，使 $\lambda \vv{OA} \; + m\vv{OB} + n\vv{OC} = \vv{0}$ ，那么 $\lambda  + m + n$ 的值为\blkda{0}.

    \item 已知空间向量 $\vv{a},\vv{b},\left| \vv{a}\right|  = 3\sqrt{2},\left| \vv{b}\right|  = 5,\vv{m} = \vv{a} + \vv{b},\vv{n} = \vv{a} + \lambda \vv{b}, < \vv{a},\vv{b} >  = {135}^{ \circ  }$ ,若 $\vv{{m}} \bot  \vv{{n}}$ , 则 $\lambda$ 的值为\blkda{$-\dfrac{3}{10}$}.

    \item 设 $\vv{a} \bot  \vv{b},\langle \vv{a},\vv{c}\rangle  = \frac{\pi }{3},\langle \vv{b},\vv{c}\rangle  = \frac{\pi }{6}$ ，且 $\left| \vv{a}\right|  = 1$ ， $\left| \vv{b}\right|  = 2$ ， $\left| \vv{c}\right|  = 3$ ，则 $\left| {\vv{a} + \vv{b} + \vv{c}}\right|  =$ \blkda{$\sqrt{17+6\sqrt{3}}$}.

    \item 在平行六面体 ${ABCD} - {EFGH}$ 中, $\vv{AG} = x\vv{AC} + y\vv{AF} + z\vv{AH}$ ,则 $x + y + z =$ \blkda{$\dfrac{3}{2}$}. 
    
    \item 在三棱柱 ${ABC} - {A}_{1}{B}_{1}{C}_{1}$ 中， $E$ 是 ${AB}$ 的中点， $F$ 是 ${C}_{1}C$ 的中点，若 $\vv{AB} = \vv{a},\vv{AC} = \vv{b}$ ， $\vv{{A}_{1}A} = \vv{c}$ ，则 $\vv{EF}$ 可表示为 \dotfill{(\kh{B})}
    \xx{$\frac{1}{2}\vv{a} + \frac{1}{2}\vv{b} - \vv{c}$ }
    {$- \frac{1}{2}\vv{a} + \vv{b} - \frac{1}{2}\vv{c}$ }
    {$\vv{a} - \frac{1}{2}\vv{b} + \frac{1}{2}\vv{c}$} 
    {$\frac{1}{2}\vv{a} + \frac{1}{2}\vv{b} - \frac{1}{2}\vv{c}$}

    
    \item 下列说法中在平面向量范围中是正确的但在空间向量中是错误的\dotfill{(\kh{D})}  \xx{$\vv{a}//\vv{b},\vv{b}//\vv{c} \Rightarrow$} {$\vv{a}//\vv{b},\vv{b}//\vv{c} \Rightarrow  \vv{a}//\vv{c}$ }{$\vv{a}//\vv{c} \vv{a} \bot  \vv{b},\vv{b} \bot  \vv{c} \Rightarrow  \vv{a}//\vv{c}$ } 
    {$\vv{a} \bot  \vv{b},\vv{b} \bot  \vv{c},\vv{b} \neq  \vv{0} \Rightarrow  \vv{a}//\vv{c}$}

    \item 若 $\vv{a},\vv{b},\vv{c}$ 为非零向量,则“存在非零实数 $a,\beta ,\gamma$ 使得 $a\vv{a} + \beta \vv{b} + \gamma \vv{c} = \vv{0}$ ”是“ $\vv{a},\vv{b},\vv{c}$ 三向量共面”的\blkda{充分非必要}条件

    \item 已知空间四边形 ${OABC}$ 中， $\angle {AOB} = \angle {BOC} = \angle {AOC}$ ，且 ${OA} = {OB} = {OC}$ 。M、N 分别是 OA、BC 的中点，G是 MN 的中点。求证： ${\text{ OG }\bot \text{ BC }}$ 。（用向量法）
    \begin{jieda}
        $\vv{OG}=\dfrac{1}{4}(\vv{OA}+\vv{OB}+\vv{OC})$,$\vv{BC}=\dfrac{1}{2}(\vv{OC}-\vv{OB})$,因此$\vv{OG} \cdot \vv{BC}=0 $.
    \end{jieda}


    \item 已知平行六面体 ${ABCD} - {A}^{\prime }{B}^{\prime }{C}^{\prime }{D}^{\prime }$ ， $M$ 是 $A{A}^{\prime }$ 的中点，点 $G$ 在对角线 ${A}^{\prime }C$ 上且 ${CG} : G{A}^{\prime } \; = 2 : 1$ ,设 $\vv{CD} = a,\vv{CB} = b,\vv{C{C}^{\prime }} = c$ ,试用 $a\text{ 、 }b\text{ 、 }c$ 表示 $\vv{CA}\text{ 、 }\vv{C{A}^{\prime }} \cdot  \vv{CM}\text{ 、 }\vv{CG}$ .
    \begin{jieda}
        $\vv{CA}=\vv{a}+\vv{b}$;\\
        $\vv{CA'}=\vv{a}+\vv{b}+\vv{c}$;\\
        $\vv{CM}=\vv{a}+\vv{b}+\dfrac{1}{2}\vv{c}$\\
        $\vv{CG}=\dfrac{2}{3}(\vv{a}+\vv{b}+\vv{c})$
    \end{jieda}

    \item 已知空间四边形 ${ABCD}$ 的每条边和对角线长都等于 $a$ ,点 $E\text{ 、 }F\text{ 、 }G$ 分别是 ${AB}\text{ 、 }{AD}$ 、 ${DC}$ 的中点,求下列向量的数量积.

(1) $\vv{AB} \cdot  \vv{AC}$ ; (2) $\vv{AD} \cdot  \vv{BD}$ ; (3) $\vv{GF} \cdot  \vv{AC}$ ; (4) $\vv{EF} \cdot  \vv{BC}$ .

\begin{jieda}
    (1)$\dfrac{a^2}{2}$;(2)$\dfrac{a^2}{2}$;(3)$-\dfrac{a^2}{2}$;(4)$\dfrac{a^2}{4}$
\end{jieda}

    \item 如图，在四棱锥 $P - {ABCD}$ 中，底面为直角梯形， ${AD}//{BC}$ ， $\angle {BAD} = {90}^{ \circ  }$ ， ${PA} \bot$ 底面 ${ABCD}$ ,且 ${PA} = {AD} = {AB} = {2BC},M,N$ 分别为 ${PC},{PB}$ 的中点. (用向量法)
    \begin{figure}[h]
        \centering
        \includegraphics[width=0.25\linewidth]{bodyxb//src/'ifjaoiuf]aouhf'ouhwqefqwe.png}
    \end{figure}
    
    \begin{enumerate}
        \item 求证： ${PB}\bot {DM}$ ；
        \item 求 ${AC}$ 与 ${PD}$ 所成角的余弦值.
    \end{enumerate}

    \begin{jieda}
        \begin{enumerate}
            \item $\vv{PB}=\vv{AB}-\vv{AP}$\\
            $\vv{DM}=\dfrac{1}{2}(\vv{AP}-\dfrac{3}{2}\vv{AD}+\vv{AB})$\\
            因此$\vv{PB}\cdot \vv{DM}=0$;
            \item 设$|AB|=t$,\\
            $\vv{AC}=\dfrac{1}{2}\vv{AD}+\vv{AB}$,\\
            $\vv{PD}=\vv{AD}-\vv{AP}$,\\
            从而$|\vv{AC}|=\dfrac{\sqrt{5}}{2}t,|\vv{PD}|=\sqrt{2}t, \vv{AC}\cdot\vv{PD}=\dfrac{1}{2}t^2$,\\
            因此$cos<\vv{AC},\vv{PD}>=\dfrac{|\vv{AC}\cdot\vv{PD}|}{|\vv{AC}||\vv{PD}|=\dfrac{\sqrt{10}}{10}}$
        \end{enumerate}
    \end{jieda}


\end{enumerate}


\section*{空间向量的基本定理}

\begin{enumerate}

\item 如下三个命题中，正确的命题是\blkda{(2)(3)}

\begin{enumerate}
\item 如果向量 $\vv{a}$ 、 $\vv{b}$ 与任何向量不能构成空间向量的一个基底，那么 $\vv{a}$ 与 $\vv{b}$ 不共线；

\item $O$ 、 $A$ 、 $B$ 、 $C$ 为空间四个点，且向量 $\vv{OA}$ 、 $\vv{OB}$ 、 $\vv{OC}$ 不构成空间的一个基底， 那么 $O\text{ 、 }A\text{ 、 }B\text{ 、 }C$ 四点必共面;

\item 已知向量 $\vv{a}\text{ 、 }\vv{b}\text{ 、 }\vv{c}$ 是空间的一个基底，则向量 $\vv{a} + \vv{b}\text{ 、 }\vv{a} - \vv{b}\text{ 、 }\vv{c}$ 也是空间的一个基底
 
\end{enumerate}


\item 空间四边形 ${O}{{ABC}}$ 中， $\vv{{O}{A}} = \vv{{a}},\vv{{O}{B}} = \vv{{b}},\vv{{O}{C}} = \vv{{c}}$ ，且 $\vv{{O}{M}} = \frac{2}{3}\vv{{O}{A}}$ ， $\vv{BN} = \vv{NC}$ ，则 $\vv{MN} =$ \dotfill{(\kh{D})}

\xx{$\frac{1}{2}\vv{{a}} - \frac{2}{3}\vv{{b}} + \frac{1}{2}\vv{{c}}$ }
{$\frac{1}{2}\vv{{a}} + \frac{1}{2}\vv{{b}} - \frac{1}{2}\vv{{c}}$}
{$- \frac{2}{3}\vv{a} + \frac{2}{3}\vv{b} + \frac{1}{2}\vv{c}$} 
{$- \frac{2}{3}\vv{a} + \frac{1}{2}\vv{b} + \frac{1}{2}\vv{c}$}


\item 已知三棱柱 ${ABC} - {A}^{\prime }{B}^{\prime }{C}^{\prime }$ 中， $M$ 是 ${B}^{\prime }C$ 和 $B{C}^{\prime }$ 的交点，设 $\vv{A{A}^{\prime }} = \vv{a}$ ， $\vv{AB} = \vv{b}$ ， $\vv{AC} = \vv{c}$ ，则 $\vv{AM} =$ \blkda{$\dfrac{1}{2}(\vv{a}+\vv{b}+\vv{c})$}.

\item 已知 $\vv{{a}}\text{ 、 }\vv{{b}}\text{ 、 }\vv{{c}}$ 为不共面的三个空间向量，若 $\vv{{m}} = \vv{{a}} - \vv{{b}} + \vv{{c}}$ 与 $\vv{{n}} = {x}\vv{{a}} + {y}\vv{{b}} + \; \vv{{c}}$ 共线，则 ${x} \neq  {y}$ 的值为\blkda{0}.

\item 已知 ${A}\left( {1,1,0}\right) ,{B}\left( {1,0, - 1}\right) ,{C}\left( {1,x + 2,{2x}}\right)$ 三点共线，则 ${x} =$\blkda{1}

\item 已知正方体 ${ABCD} - {A}_{1}{B}_{1}{C}_{1}{D}_{1}$ ， $E$ ， $M$ ， $N$ 分别是棱 ${D}_{1}D,{A}_{1}{D}_{1},{D}_{1}{C}_{1}$ 的中点， $O$ 是线段 ${BD}$ 的中点， $\vv{{B}_{1}O},\vv{EM},\vv{EN}$ 是否共面\blkda{是}。

\item 已知 $\vv{a},\vv{b},\vv{c}$ 不共面，且 $\vv{d} = 2\vv{a} + x\vv{b} - 3\vv{c},\vv{e} = y\vv{a} - {12}\vv{b} + 9\vv{c},\vv{d}//\vv{e}$ ，则实数 $x = \blkda{4},y =$ \blkda{-6}。


\item 已知正方体 ${ABCD} - {A}_{1}{B}_{1}{C}_{1}{D}_{1}$ ，点 $E$ 是棱 ${BC}$ 的中点， ${{D}_{1}B}$ 与 ${A}_{1}E$ 相交于点 $F$ ， 设 $\vv{{D}_{1}{A}_{1}} = \vv{a},\vv{{D}_{1}D} = \vv{b},\vv{{D}_{1}{C}_{1}} = \vv{c}$ ，则 $\vv{DF} =$ \blkda[4]{$\dfrac{2}{3}\vv{a}-\dfrac{1}{3}\vv{b}+\dfrac{2}{3}\vv{c}$}。

\item 已知 $\vv{a},\vv{b},\vv{c}$ 不共面， $\vv{{d}_{1}} = 2\vv{a} + 3\vv{b} - \vv{c},\vv{{d}_{2}} = \vv{a} - 4\vv{b} + 2\vv{c},\vv{d} = \vv{{d}_{1}} + m\vv{{d}_{2}},\vv{d}$ 与 $\vv{b},\vv{c}$ 共面， 则实数 $m =$ \blkda{-2}。

\item 已知向量 $\vv{a},\vv{b},\vv{c}$ 的模分别为1,2,3，且它们两两的夹角都是 $\frac{\pi }{3}$ ， $\left| {3\vv{a} + \frac{3}{2}\vv{b} + \vv{c}}\right|  =$\blkda{$3\sqrt{6}$}

\item 已知向量 $\vv{a}\text{ 、 }\vv{b}$ ，且 $\vv{AB} = \vv{a} + 2\vv{b},\vv{BC} =  - 5\vv{a} + 6\vv{b},\vv{CD} = 7\vv{a} - 2\vv{b}$ ，则一定共线的三点是\dotfill{(\kh{A})}

\xx{$A$ 、 $B$ 、 $D$} {$A$ 、 $B$ 、 $C$} { $B$ 、 $C$ 、 $D$} {$A$ 、 $C$ 、 $D$}

\item  $A,B,C$ 不共线，对空间任意一点 $O$ ，若 $\vv{OP} = \frac{3}{4}\vv{OA} + \frac{1}{8}\vv{OB} + \frac{1}{8}\vv{OC}$ ，则 $P,A$ ，
$B,C$ 四点\dotfill{(\kh{B})}

\xx{不共面} {共面}  {不一定共面}{无法判断是否共面}

\item \在四面体 ${ABCD}$ 中，“ ${AB}\bot {CD}$ ” 是 “ ${A{C}^{2}} + {B{D}^{2}} = {A{D}^{2}} + {B{C}^{2}}$ ” 的\blkda{充要}条件。

\item 在平行四边形 ${ABCD}$ 中， ${AB} = {AC} = 1,\angle {ACD} = {90}^{ \circ  }$ ，将它沿对角线 ${AC}$ 折起， 使 ${AB}$ 与 ${CD}$ 成 ${60}^{ \circ  }$ ，则 $B,D$ 的距离是\blkda{2或$\sqrt{2}$}。

\item 在空间四边形 ${ABCD}$ 中，下列结论正确的是\dotfill{(\kh{A})}

\xx{$\vv{AB} \cdot  \vv{CD} + \vv{AC} \cdot  \vv{DB} + \vv{AD} \cdot  \vv{BC} = 0$ }
{$\vv{AB} \cdot  \vv{CD} + \vv{AC} \cdot  \vv{DB} = \vv{AD} \cdot  \vv{BC}$}
{$\vv{AB} \cdot  \vv{CD} = \vv{AC} \cdot  \vv{DB} + \vv{AD} \cdot  \vv{BC}$ }
{以上都不对}

\item 设 $A,B,C,D$ 是 空 间 不 共 面 的 四 点 ，且 满 足 $\vv{AB} \cdot  \vv{AC} = 0,\vv{AB} \cdot  \vv{AD} = 0,\vv{AC} \cdot  \vv{AD} = 0$ ,则 $\bigtriangleup  {BCD}$ 是 \dotfill{(\kh{C})}

\xx{钝角三角形}{直角三角形}{锐角三角形}{不能确定}


\item 若 $\vv{a},\vv{b},\vv{c}$ 为非零向量，则 “ $\vv{a}\bot \vv{b},\vv{a}\bot \vv{c}$ ” 是 “ $\vv{a}\bot \left( {\vv{b} + \vv{c}}\right)$ ” 的\blkda[3]{充分非必要条件}条件。


\item 已知空间四边形 ${ABCD}$ ,连结 ${AC},{BD}$ ,设 $M,G$ 分别是 ${BC},{CD}$ 的中点,化简下列各表达式:
\begin{enumerate}
    \item $\vv{AB} + \vv{BC} + \vv{CD}$ ；
    \item $\vv{AB} + \frac{1}{2}\left( {\vv{BD} + \vv{BC}}\right)$ ； 
    \item $\vv{AG} - \frac{1}{2}\left( {\vv{AB} + \vv{AC}}\right)$ 。
\end{enumerate}
\begin{jieda}
    (1)$\vv{AD}$;(2)$\vv{AG}$;(3)$\vv{MG}$.
\end{jieda}


\item 已知平行四边形 ${ABCD}$ ，从平面 ${AC}$ 外一点 $O$ 引向量。 $\vv{OE} = k\vv{OA},\vv{OF} = k\vv{OB},\vv{OG} = k\vv{OC},\vv{OH} = k\vv{OD}$ 。
\begin{enumerate}
    \item 求证：四点 $E,F,G,H$ 共面；
    \item 平面 ${AC}//$ 平面 ${EG}$ 。
\end{enumerate}

\begin{jieda}
    \begin{enumerate}
        \item $ \because $ 四边形 $ {ABCD} $ 是平行四边形, $ \therefore \vv{AC} = \vv{AB} + \vv{AD} $ ,
$ \because \vv{EG} = \vv{OG} - \vv{OE} $ ,
$ = k \cdot  \vv{OC} - k \cdot  \vv{OA} = k\left( {\vv{OC} - \vv{OA}}\right)  = k\vv{AC} = k\left( {\vv{AB} + \vv{AD}}\right) $
$ = k\left( {\vv{OB} - \vv{OA} + \vv{OD} - \vv{OA}}\right)  = \vv{OF} - \vv{OE} + \vv{OH} - \vv{OE} = \vv{EF} + \vv{EH} $
$ \therefore E $ 、 $ F $ 、 $ G $ 、 $ H $ 四点共面;
    \item $ \because \vv{EF} = \vv{OF} - \vv{OE} = k\left( {\vv{OB} - \vv{OA}}\right)  = k \cdot  \vv{AB},\therefore {EF}\parallel {AB} $
又因为 $ {EF} ⊄ $ 平面 $ {ABCD} $ , $ {AB} \subset $ 平面 $ {ABCD} $ ,所以 $ {EF}// $ 平面 $ {ABCD} $
又 $ \because \vv{EG} = k \cdot  \vv{AC},\therefore {EG}\parallel {AC} $ ,
$ {EG} ⊄ $ 平面 $ {ABCD},{AC} \subset $ 平面 $ {ABCD},{EG} \parallel $ 平面 $ {ABCD} $ ,
又 $ {EF} \cap  {EG} = E,{EF},{EG} \subset $ 平面 $ {EG} $
所以,平面 $ {EG} \parallel $ 平面 $ {AC} $
    \end{enumerate}.
    
\end{jieda}



\item 已知平行六面体 ${ABCD} - {A}^{\prime }{B}^{\prime }{C}^{\prime }{D}^{\prime }$ 中， ${AB} = 4,{AD} = 3,{AA}^{\prime } = 5,\angle {BAD} = {90}^{ \circ  }$ ， $\angle {BA}{A}^{\prime } = \angle {DA}{A}^{\prime } = {60}^{ \circ  }$ ,求 $A{C}^{\prime }$ 的长。

\begin{jieda}
可得 $ \vv{A{C}^{\prime }} = \vv{AC} + \vv{C{C}^{\prime }} = \vv{AB} + \vv{AD} + \vv{A{A}^{\prime }} $ ,
$ {\left| \vv{AC}\right| }^{2} = {\left| \vv{AB} + \vv{AD} + \vv{A{A}^{\prime }}\right| }^{2} = {\vv{AB}}^{2} + {\vv{AD}}^{2} + {\vv{AA'}}^{2}+2\left( {\vv{AB} \cdot  \vv{AD} + \vv{AB} \cdot  \vv{A{A}^{\prime }} + \vv{AD} \cdot  \vv{A{A}^{\prime }}}\right) $
$ = {4}^{2} + {3}^{2} + {5}^{2} + 2\;\left( {4 \times  3 \times  0 + 4 \times  5 \times  \frac{1}{2} + 3 \times  5 \times  \frac{1}{2}}\right)  = {85} $,
故 $ {AC} $ 的长等于 $ \left| \vv{AC}\right|  = \sqrt{85} $.
\end{jieda}

\item 我们知道: 在平面中任意不共线三点 $O$ 和 $A,B$ ,若 $\vv{OP} = x\vv{OA} + y\vv{OB}\left( {x + y = 1}\right)$ , 则 $P,A,B$ 三点共线。对空间任一点 $O$ 和不共线的三点 $A,B,C$ ,问满足向量式: $\vv{OP} = x\vv{OA} + y\vv{OB} + z\vv{OC}$ (其中 $x + y + z = 1$ ) 的四点 $P,A,B,C$ 的位置有何关系?
\begin{jieda}
    $\vv{OP}=x\vv{OA}+y\vv{OB}+(1-x-y)\vv{OC}$,从而$\vv{CP}=x\vv{CA}+y\vv{CB}$,三个向量共面，所以四点共面。
\end{jieda}

\end{enumerate}




\section*{空间向量的坐标表示}

A组:
\begin{enumerate}

\item 判断下列命题的真假（正确的打 “✓”，错误的打 “✗”）
\begin{enumerate}

    \item 以原点为始点的向量 $\vv{OP}$ 的坐标和点 $P$ 的坐标相同；（ \blkda{\ding{51}}）

    \item 在空间直角坐标系中，在 ${Ox}$ 轴上的点一定是 $\left( {0,b,c}\right)$ ；（ \blkda{\ding{53}}）

    \item 在空间直角坐标系中，在 ${xOz}$ 平面上的点的坐标为 $\left( {a,0,c}\right)$ ； \blkda{\ding{51}}）

    \item 空间直角坐标系中 ${xOz}$ 平面上点的坐标满足 $z = 0;$ （ \blkda{\ding{53}}）

    \item 关于坐标平面 ${yOz}$ 对称的点的坐标其纵、竖坐标不变，横坐标相反；（ \blkda{\ding{51}}）

\end{enumerate}

    \item 设 $\left\{  {\vv{i},\vv{j},\vv{k}}\right\}$ 是空间向量的一个单位正交基底,则向量 $\vv{a} = 3\vv{i} + 2\vv{j} - \vv{k},\vv{b} =  - 2\vv{i} + 4\vv{j} + 2\vv{k}$ 的坐标分别是\blkda[4]{$(3,2,-1)$和$(-2,4,2)$}.

    \item 已知正方体 ${ABCD} - {A}_{1}{B}_{1}{C}_{1}{D}_{1}$ 的棱长为 1,若以 $\{ \vv{AB},\vv{AD},\vv{A{A}_{1}}\}$ 为基底，则 $\vv{A{C}_{1}} \; =$ \blkda[4]{$\vv{AB}+\vv{AD}+\vv{AA_1}$}. $\vv{A{C}_{1}}$ 的坐标是\blkda{$(1,1,1)$}.

    \item 向量 $\vv{a} = \left( {0,1,0}\right) ,\vv{b} = \left( {1, - 1,0}\right)$ 的夹角大小为\blkda{$\dfrac{3\pi}{4}$}.

    \item 已知向量 $\vv{a} = \left( {-m,1,3}\right) ,\vv{b} = \left( {2,n,1}\right)$ 平行，则 $m =$ \blkda{-6}， $n =$\blkda{$\dfrac{1}{3}$}。

    \item 若 $\{ \vv{a},\vv{b},\vv{c}\}$ 是空间的一个基底,试判断 $\{ \vv{a} + \vv{b},\vv{b} + \vv{c},\vv{c} + \vv{a}\}$ 能否作为该空间的一个基底.
    \begin{jieda}
        因为$\vv{a}+\vv{b}=x(\vv{b}+\vv{c})+y(\vv{a}+\vv{c})$无解，因此三个向量不共面，因此可以作为空间中的一组基底。
    \end{jieda}

\end{enumerate}


B 组:
\begin{enumerate}

    \item 在三角形 ${{ABC}}$ 中，已知 $\vv{AB} = \left( {2,4,0}\right) ,\vv{BC} = \left( {-1,3,0}\right)$ ，则 $\angle {ABC} =$ \blkda{$\dfrac{3\pi}{4}$}.

    \item 已知空间的一个基底 $\{ \vv{a},\vv{b},\vv{c}\} ,\vv{m} = \vv{a} - \vv{b} + \vv{c},\vv{n} = x\vv{a} + y\vv{b} + \vv{c}$ ,若 $\vv{m}$ 与 $\vv{n}$ 共线，则 $x + y$ 等于\blkda{0}.


    \item 在直三棱柱 ${ABO} - {A}_{1}{B}_{1}{O}_{1}$ 中， $\angle {AOB} = \frac{\pi }{2},{AO} = 4,{BO} = 2,{A{A}_{1}} = 4,D$ 为 ${A}_{1}{B}_{1}$的中点，则在如图所示的空间直角坐标系中， $\vv{DO}$ 的坐标是\blkda[2]{$(-2,-1,-4)$}； $\vv{{A}_{1}B}$ 的坐标是\blkda[2]{$(-4,2,-4)$}.

\begin{figure}[h]
    \centering
    \includegraphics[width=0.25\linewidth]{bodyxb//src/seoihgaoijhrfaifra.png}
\end{figure}

    \item 如图，在四面体 ${ABCD}$ 中， $G$ 为 $\bigtriangleup  {ABC}$ 的重心， $E$ 是 ${BD}$ 上一点， ${BE} = {3ED}$ ，以 $\{ \vv{AB},\vv{AC}$ ， $\vv{AD}\}$ 为基底，则 $\vv{GE} =$ \blkda[4]{$-\dfrac{1}{12}\vv{AB}-\dfrac{1}{3}\vv{AC}+\dfrac{3}{4}\vv{AD}$}.

\begin{figure}[h]
    \centering
    \includegraphics[width=0.25\linewidth]{bodyxb//src/soghaoeuhrfaoewfa.png}
\end{figure}

    \item 已知点 $A$ 在基底 $\{ \vv{a},\vv{b},\vv{c}\}$ 下的坐标为 $\left( {2,1,3}\right)$ ，其中 $\vv{a} = 4\vv{i} + 2\vv{j},\vv{b} = 2\vv{j} \; + 3\vv{k}$ ， $\vv{c} = {3\vv{k}} - \vv{j}$ ，则点 $A$ 在基底 $\{ \vv{i}$ ， $\vv{j}$ ， $\vv{k}\}$ 下的坐标为\blkda{$(8,3,12)$}.


    \item 给出下列命题，其中正确命题的序号为\blkda[3]{(1)(2)(3)(4)}.

\begin{enumerate}
    \item 空间任意三个不共面的向量都可以作为一个基底；
    \item 已知向量 $\vv{a}//\vv{b}$ ，则 $\vv{a}$ ， $\vv{b}$ 与任何向量都不能构成空间的一个基底；
    \item $A$ 、 $B$ 、 $M$ 、 $N$ 是空间四点，若 $\vv{BA}$ 、 $\vv{BM}$ 、 $\vv{BN}$ 不能构成空间的一个基底，那么 $A$ 、 $B\text{ 、 }M\text{ 、 }N$ 共面;
    \item 已知向量组 $\{ \vv{a},\vv{b},\vv{c}\}$ 是空间的一个基底，若 $\vv{m} = \vv{a} + \vv{c}$ ，则 $\{ \vv{a},\vv{b},\vv{m}\}$ 也是空间的一个基底.
    
\end{enumerate}

    
    \item 如图，四个棱长为 1 的正方体排成一个正四棱柱， ${AB}$ 是一条侧棱, $\;{P}_{i}\left( {i = 1,2,\cdots ,8}\right)$ 是上底面上其余的八个点,则 $\vv{AB} \cdot  \vv{A{P}_{i}}\left( {i = 1,2,\cdots ,8}\right)$ 的不同值的个数为\blkda{1}.

\begin{figure}[h]
    \centering
    \includegraphics[width=0.25\linewidth]{bodyxb//src/fakbhf'aohfouhwfa.png}
\end{figure}

    \item 如图，在四棱锥 $P = {ACD}$ 中，底面 ${ABCD}$ 是正方形， $E$ 为 ${PD}$ 中点， 若 $\vv{PA} = \vv{a},\vv{PB} = \vv{b},\vv{PC} = \vv{c}$ ，则 $\vv{BE} =$ \dotfill{(\kh{C})}
    
    \xx{$\frac{1}{2}\vv{a} - \frac{1}{2}\vv{b} + \frac{1}{2}\vv{c}$}
    
    {$\frac{1}{2}\vv{a} - \frac{1}{2}\vv{b} - \frac{1}{2}\vv{c}$}
    
    {$\frac{1}{2}\vv{a} - \frac{3}{2}\vv{b} + \frac{1}{2}\vv{c}$ }
    {$\frac{1}{2}\vv{a} - \frac{1}{2}\vv{b} + \frac{3}{2}\vv{c}$}

\begin{figure}[h]
    \centering
    \includegraphics[width=0.25\linewidth]{bodyxb//src/sgoidhgoauehrfoawrfa.png}
\end{figure}

    \item 已知正方体 ${ABCD} - {A}_{1}{B}_{1}{C}_{1}{D}_{1}$ 中， $\vv{{A}_{1}E} = \frac{1}{4}\vv{{A}_{1}{C}_{1}}$ ，若 $\vv{AE} = x\vv{A{A}_{1}} + y\left( {\vv{AB} + \vv{AD}}\right)$ ，则 \dotfill{(\kh{D})}
    \xx{$x = 1,y = \frac{1}{2}$}
    {$x = \frac{1}{2},y = 1$ }
    {$x = 1,y = \frac{1}{3}$}
    {$x = 1,y = \frac{1}{4}$}

    \item 在空间直角坐标系中,点 $P\left( {-2,1,4}\right)$ . (1) 求点 $P$ 关于 $x$ 轴的对称点的坐标；（2）求点 $P$ 关于 ${xOy}$ 平面的对称点的坐标；（3）求点 $P$ 关于点 $M\left( {2,}\right.$ -4)的对称点的坐标;
    \begin{jieda}
        (1)$(-2,-1,-4)$;(2)$(-2,1,-4)$;(3)$(6,-3,-12)$.
    \end{jieda}


    \item 如图所示，PA 垂直于正方形 ${ABCD}$ 所在的平面， $M,N$ 分别是 ${AB},{PC}$ 的中点,并且 ${PA} = {AB} = 1$ ; 试建立适当的空间直角坐标系并求：向量 $\vv{MN}$ 的坐标.
\begin{figure}[h]
    \centering
    \includegraphics[width=0.25\linewidth]{bodyxb//src/ihfjaouhfaouwhf.png}
\end{figure}

    \begin{jieda}
        以$\vv{AB}$为x轴正方向，以$\vv{AD}$为y轴正方向，以$\vv{AP}$为z轴正方向建立空间直角坐标系。$\vv{MN}=(0,\dfrac{1}{2},\dfrac{1}{2})$
    \end{jieda}


    \item 已知 $\vv{a} = \vv{i} + \vv{j} + \vv{k},\vv{b} = \vv{i} - 2\vv{j} + 2\vv{k}$ ,求:
    \begin{enumerate}
        \item 向量 $\vv{a}$ 与 $\vv{b}$ 的夹角的大小;
        \item 向量 $\vv{a}$ 与 $\vv{b}$ 所在的直线的夹角大小
    \end{enumerate}
    \begin{jieda}
        (1)$arccos\dfrac{\sqrt{3}}{9}$;(2)$arccos\dfrac{\sqrt{3}}{9}$;
    \end{jieda}
    
    \item 已知向量 $\vv{a} = \left( {1, - 3,2}\right) ,\vv{b} = \left( {2,0, - 8}\right)$ ，求单位向量 $\vv{c}$ ，使 $\vv{c}$ 与向量 $\vv{a}$ 、 $\vv{b}$ 都垂直.
    \begin{jieda}
        $(\dfrac{4\sqrt{21}}{21},\dfrac{2\sqrt{21}}{21},\dfrac{\sqrt{21}}{21})$和$(-\dfrac{4\sqrt{21}}{21},-\dfrac{2\sqrt{21}}{21},-\dfrac{\sqrt{21}}{21})$,
    \end{jieda}


    \item 在长方体 ${ABCD} - {A}^{\prime }{B}^{\prime }{C}^{\prime }{D}^{\prime }$ 中,已知 ${AB} = 6,{AD} = 2,A{A}^{\prime } = 1,P$ 是 ${AB}$ 上的点且 ${PB} = {2AP},M$ 是 ${DC}$ 上的点，且 ${DM} = {2MC},N$ 是 ${B}^{\prime }{C}^{\prime }$ 的中点，求直线 $P{D}^{\prime }$ 与 ${MN}$ 所成的角6的大小.
    \begin{jieda}
        $arccos\dfrac{5\sqrt{6}}{18}$
    \end{jieda}

    
    \item 已知正三棱柱 ${ABC} - {A}_{1}{B}_{1}{C}_{1},{AB} = a,{A{A}_{1}} = b$ ，间：当 $a$ 与 $b$ 满足什么条件时，侧面对角线 $A{B}_{1}$ 垂直于 $B{C}_{1}$ .
    \begin{jieda}
        首先,建立空间直角坐标系。设正三棱柱底面 $ {ABC} $
在 $ {xOy} $ 平面上, $ A\left( {0,0,0}\right) $ , $ B\left( {a,0,0}\right) $ ,
$ C\left( {\frac{a}{2},\frac{\sqrt{3}}{2}a,0}\right) $ ,上底面顶点坐标为 $ {A}_{1}\left( {0,0,b}\right) $ ,
$ {B}_{1}\left( {a,0,b}\right) ,{C}_{1}\left( {\frac{a}{2},\frac{\sqrt{3}}{2}a,b}\right) $ 。
侧面对角线向量为:
$ \vv{A{B}_{1}} = {B}_{1} - A = \left( {a,0,b}\right) ,\;\vv{B{C}_{1}} = {C}_{1} - B = \left( {-\frac{a}{2}}\right. $ ,
两直线垂直的充要条件是其方向向量点积为零:
$ \overset{\rightarrow }{A{B}_{1}} \cdot  \overset{\rightarrow }{B{C}_{1}} = a \cdot  \left( { - \frac{a}{2}}\right)  + 0 \cdot  \frac{\sqrt{3}}{2}a + b \cdot  b =  - \frac{{a}^{2}}{2} + {b}^{2} $
解得：
$a = \sqrt{2}b $ .
    \end{jieda}

    \item 正四棱柱 ${ABCD} - {A}_{1}{B}_{1}{C}_{1}{D}_{1}$ 中，底面 ${ABCD}$ 是边长为 4 的正方形， $A{C}_{1}$ 与 ${BDA}$ 交于点 $N,B{C}_{1}$ 与 ${B}_{1}C$ 交于点 $M$ ，且 ${AM}\bot {BN}$ . （1）用向量方法求 $A{A}_{1}$ 的长. （2） 对于 $r$ 个向量 $\vv{{a}_{1}}\text{ 、 }\vv{{a}_{2}}\text{ 、 }\cdots \text{ 、 }\vv{{a}_{n}}$ ,如果存在不全为零的 $r$ 个实数 ${\lambda }_{1}\text{ 、 }{\lambda }_{2}\text{ 、 }\cdots \text{ 、 }{\lambda }_{n}$ , 使得 ${\lambda }_{1}\vv{{a}_{1}} + {\lambda }_{2}\vv{{a}_{2}} + \cdots  + {\lambda }_{n}\vv{{a}_{n}} = \vv{0}$ ，则称这 $n$ 个向量 $\vv{{a}_{1}}\text{ 、 }\vv{{a}_{2}}\text{ 、 }\cdots \text{ 、 }\vv{{a}_{n}}$ 线性相关，否则称为线性无关.试判断 $\vv{AM}\text{ 、 }\vv{BN}\text{ 、 }\vv{CN}$ 是否线性相关.

    \begin{jieda}
        \begin{enumerate}
            \item 由题设可得 $ \vv{AM} = \vv{AB} + \frac{1}{2}\vv{B{B}_{1}} + \frac{1}{2}\vv{BC} = \vv{AB} + \frac{1}{2}\vv{A{A}_{1}} + \frac{1}{2}\vv{AD} $ ,
$ \vv{BN} = \vv{B{B}_{1}} + \frac{1}{2}\vv{{B}_{1}{A}_{1}} + \frac{1}{2}\vv{{B}_{1}{C}_{1}} =  - \frac{1}{2}\vv{AB} + \frac{1}{2}\vv{AD} + \vv{A{A}_{1}}, $
故 $ \left( {\vv{AB} + \frac{1}{2}\vv{A{A}_{1}} + \frac{1}{2}\vv{AD}}\right)  \cdot  \left( {-\frac{1}{2}\vv{AB} + \frac{1}{2}\vv{AD} + \vv{A{A}_{1}}}\right)  = 0 $ ,
整理得到 $ - \frac{1}{2} \times  {16} + \frac{1}{4} \times  {16} + \frac{1}{2}{\vv{A{A}_{1}}}^{2} = 0 $ ,故 $ \left| \vv{A{A}_{1}}\right|  = 2\sqrt{2} $ .
    \item 令 $ {\lambda }_{1}\vv{AM} + {\lambda }_{2}\vv{BN} + {\lambda }_{3}\vv{CD} = \vv{0} $ ,
则 $ {\lambda }_{1}\left( {\vv{AB} + \frac{1}{2}\vv{A{A}_{1}} + \frac{1}{2}\vv{AD}}\right)  + {\lambda }_{2}\left( {-\frac{1}{2}\vv{AB} + \frac{1}{2}\vv{AD} + \vv{A{A}_{1}}}\right)  + {\lambda }_{3}\vv{AB} = \vv{0} $ ,
整理得到 $ \left( {{\lambda }_{1} - \frac{1}{2}{\lambda }_{2} + {\lambda }_{3}}\right) \vv{AB} + \left( {\frac{1}{2}{\lambda }_{1} + {\lambda }_{2}}\right) \vv{A{A}_{1}} + \left( {\frac{1}{2}{\lambda }_{1} + \frac{1}{2}{\lambda }_{2}}\right) \vv{AD} = \vv{0} $ ,
故 $ \left\{  \begin{array}{l} {\lambda }_{1} - \frac{1}{2}{\lambda }_{2} + {\lambda }_{3} = 0 \\  \frac{1}{2}{\lambda }_{1} + {\lambda }_{2} = 0 \\  \frac{1}{2}{\lambda }_{1} + \frac{1}{2}{\lambda }_{2} = 0 \end{array}\right. $ ,解得 $ {\lambda }_{1} = {\lambda }_{2} = {\lambda }_{3} = 0 $ ,
故 $ \vv{AM},\vv{BN},\vv{CD} $ 线性无关,
        \end{enumerate}
    \end{jieda}

\end{enumerate}


\section*{空间向量在立体几何中的应用}


\subsection*{空间向量在立体几何中的应用(1)}

\begin{enumerate}
    \item 已知平面 $a$ 经过点 $A\left( {3,1, - 1}\right)$ 、 $B\left( {1, - 1,0}\right)$ ，且平行于向量 $\vv{a} = \left( {-1,0,2}\right)$ ，求平面 $a$ 的一个法向量是\blkda{$(2,-\dfrac{3}{2},1)$}.

    \item 已知点 $A\left( {5,1,3}\right) \text{ 、 }B\left( {1,6,2}\right) \text{ 、 }C\left( {5,0,4}\right) \text{ 、 }D\left( {4,0,6}\right)$ ,求过 $A\text{ 、 }B$ 且平行于 ${CD}$ 的平面的一个法向量是\blkda{$(10,9,5)$}.

    \item 在空间直角坐标系中, $\vv{{d}} = \left( {1,2, - 1}\right)$ 是直线 $l$ 的一个方向向量, $\vv{{n}} = \; \left( {2,t, - 2}\right)$ 是平面 $\alpha$ 的一个法向量. 若 $l \bot  \alpha$ ，则 $t =$ \blkda{4}.

    \item 已知空间向量 $\vv{{a}} = \left( {{x},{1}, - {2}}\right)$ 与 $\vv{{b}} = \left( {{1}, - {1},{2}}\right)$ 夹角为钝角，则实数 ${x}$ 的取值范围为\blkda[4]{$(-\infty,-1)\cup (-1,5)$}.

    \item 正四棱锥 $P - {ABCD}$ 的侧面 ${PAB}$ 是等边三角形， $E$ 为 ${PC}$ 的中点，则异面直线 ${BE}$ 和 ${PA}$ 所成角的余弦值为\blkda{$\dfrac{\sqrt{3}}{3}$}.

    \item 已知 ${ABCD} - {{A}_{1}{B}_{1}{C}_{1}{D}_{1}}$ 是长方体， ${A{A}_{1}} = 2,{AB} = {BC} = 1$ ， $E$ 为 ${BC}$ 的中点，则异面直线 ${A}_{1}E$ 与 ${D}_{1}A$ 所成角的正切值为\blkda{$\dfrac{2\sqrt{14}}{7}$}.

\begin{figure}[h]
    \centering
    \includegraphics[width=0.25\linewidth]{bodyxb//src/aushfoauhfouahefaw.png}
\end{figure}

    \item 如图，在长方体 ${ABCD} - {A}_{1}{B}_{1}{C}_{1}{D}_{1}$ 中，已知 ${AB} = 2$ ， ${AD} = A{A}_{1} = 1$ . 动点 $P$ 从 ${A}_{1}$ 出发，在棱 ${A}_{1}{B}_{1}$ 上匀速运动； 动点 $Q$ 同时从 $B$ 出发，在棱 ${BC}$ 上匀速运动， $P$ 的运动速度是 $Q$ 的两倍，各自运动到另一端点停止. 它们在运动过程中，设直线 ${PQ}$ 与平面 ${ABCD}$ 所成的角为 $\theta$ ，则 $\tan \theta$ 的取值范围是\blkda[3]{$[\dfrac{1}{2},\dfrac{\sqrt{5}}{2}]$}.

\begin{figure}[h]
    \centering
    \includegraphics[width=0.25\linewidth]{bodyxb//src/gs'uhf'aouerhFOQUaerhf.png}
\end{figure}

    \item 已知空间向量 $\vv{{O}_{{B}_{1}}}$ 、 $\vv{{O}_{{B}_{2}}}$ 、 $\vv{{O}_{{B}_{3}}}$ 两两垂直，空间中点 ${A}$ 满足 $\left| \vv{{B}_{1}}\right|  = \; \left| {\vv{AB}}_{2}\right|  = \left| {\vv{AB}}_{3}\right|  = 1$ ,记 $\vv{OP} = {\vv{OB}}_{1} + {\vv{OB}}_{2} + {\vv{OB}}_{3}$ ,则 $\left| \vv{AP}\right|$ 的取值范围为\blkda[2]{[0,\sqrt{3}]}.

    \item 已知向量 $\vv{{m}} = \left( {{a},{b},{0}}\right) ,\vv{{n}} = \left( {{c},{d},{1}}\right)$ ，其中 ${{a}}^{2} + {{b}}^{2} = {{c}}^{2} + {{d}}^{2} = 1$ ，现有以下命题:
    \begin{enumerate}
        \item 向量 $\vv{{n}}$ 与 $z$ 轴正方向的夹角恒为定值(即与 $c,d$ 无关);
        \item $\vv{{m}} \cdot  \vv{{n}}$ 的最大值为 $\sqrt{{2}}$ ；
        \item $\langle \vv{{m}},\vv{{n}}\rangle \langle \vv{{m}},\vv{{n}}$ 的夹角) 的最大值为 $\frac{3\pi }{4}$ ;
        \item 若定义 $\vv{{u}} \times  \vv{{v}} = \left| \vv{{u}}\right|  \cdot  \left| \vv{{v}}\right| \sin \langle \vv{{u}},\vv{{v}}\rangle$ ，则 $\left| {\vv{{m}} \times  \vv{{n}}}\right|$ 的最大值为 $\sqrt{2}$ .
      \end{enumerate}
    其中正确的命题有\blkda{(1)(3)(4)}. (写出所有正确命题的序号)

    \item 已知直线 $l$ 与平面 $a$ ,则 “平面 $a$ 的法向量 $\vv{n}$ 是直线 $l$ 的方向向量” 是 “ $l \bot  a$ ” 的（ \blkda{充分必要}）条件


    \item 如图，正方体 ${ABCD} - {A}^{\prime }{B}^{\prime }{C}^{\prime }{D}^{\prime }$ 的棱长为 8，点 $P$ 在左侧面 ${AD}{D}^{\prime }{A}^{\prime }$ 上，点 $P$ 到 $A{A}^{\prime }$ 的距离为 2，点 $P$ 到 ${AD}$ 的距离为 4,问在底面 ${A}^{\prime }{B}^{\prime }{C}^{\prime }{D}^{\prime }$ 上是否存在点 $Q$ ,使得 $\vv{PQ}$ 是直线 $A{C}^{\prime }$ 的方向向量? 请说明理由.注意：必须说明空间直角坐标系如何建立。

\begin{figure}[h]
    \centering
    \includegraphics[width=0.25\linewidth]{bodyxb//src/asuhf'aouwhefoauwherfohw.png}
\end{figure}

    \begin{jieda}
        存在,理由如下：
如图所示,以 $ {D}^{\prime } $ 为原点, $ {D}^{\prime }{A}^{\prime } $ 为 $ x $ 轴正方向,
$ {D}^{\prime }{C}^{\prime } $ 为 $ y $ 轴正方向, $ {D}^{\prime }D $ 为z轴正方形,建立空间直角坐标系.
则 $ A\left( {8,0,8}\right) ,{C}^{\prime }\left( {0,8,0}\right) \;P\left( {6,0,4}\right) $.
设 $ Q\left( {x,y,0}\right) $ .
$ \therefore \vv{PQ} = \left( {x - 6,y, - 4}\right) $
$ A{C}^{\prime } $ 的单位向量为 $ A\left( {-1,1, - 1}\right) $
若 $ \vv{AC} $ 是直线 $ A{C}^{\prime } $ 的方向而是
则有 $ \vv{PQ} = \lambda \vv{a} $，
从而$ \left( {x - 6,y, - 4}\right)  = \lambda \left( {-1,1, - 1}\right) $
解得点 $ Q $ 为(2,4,0)，
即存在点 $ Q $ ,使得 $ \vv{Q} $ 是直线 $ A{C}^{\prime } $ 的方向向量.
    \end{jieda}



    \item  如图，在长方体 ${ABCD} - {A}_{1}{B}_{1}{C}_{1}{D}_{1}$ 中， $E$ 、 $F$ 分别是棱 ${BC}$ ，$C{C}_{1}$ 上的点， ${CF} = {AB} = {2CE},{AB} : {AD} : {A{A}_{1}} = 1 : 2 : 4$ .
    \begin{enumerate}
        \item 求异面直线 ${EF}$ 与 ${A}_{1}D$ 所成角的大小；
        \item 证明： ${AF}\bot$ 平面 ${A}_{1}{ED}$ ；
        \item 求二面角 ${A}_{1} - {ED} - F$ 的大小.
    \end{enumerate}

\begin{figure}
    \centering
    \includegraphics[width=0.25\linewidth]{bodyxb//src/gos'ihdgousheoruhage.png}
\end{figure}

    \begin{jieda}
    以AB为x轴正方向，AD为y轴正方向，$AA'$为z轴正方向建立空间直角坐标系，设 $ {AB} = 1 $ ,依题意得 $ D\left( {0,2,0}\right) $ ,
$ F\left( {1,2,1}\right) ,{A}_{1}\left( {0,0,4}\right) ,E\left( {1,\frac{3}{2},0}\right) $
        \begin{enumerate}
            \item 易得 $ \vv{EF} = \left( {0,\frac{1}{2},1}\right) ,\vv{{A}_{1}D} = \left( {0,2, - 4}\right) $于是 $ \cos \left\langle  {\vv{EF},\vv{{A}_{1}D}}\right\rangle   = \frac{\vv{EF} ·\vv{{A}_{1}D}}{\left| \vv{EF}\right| \left| \vv{{A}_{1}D}\right| } =  - \frac{3}{5} $所以异面直线 $ {EF} $ 与 $ {A}_{1}D $ 所成角为 $ arccos\frac{3}{5} $
            \item 证明：已知 $ \vv{AF} = \left( {1,2,1}\right) ,\vv{EA} = \left( {-1, - \frac{3}{2},4}\right) ,\vv{ED} = \left( {-1,\frac{1}{2},0}\right) $于是 $ \vv{AF} \cdot  \vv{E{A}_{1}} = 0,\vv{AF} \cdot  \vv{ED} = 0 $ .因此, $ {AF} \bot  {EA},{AF} \bot  {ED} $ ,又 $ E{A}_{1} \cap  {ED} = E $所以 $ {AF} \bot $ 平面 $ A_1ED $
            \item 二面角的大小为$arccos\dfrac{2}{3}$
        \end{enumerate}
    \end{jieda}

    \item 如图，在正三棱柱 ${ABC} - {A}_{1}{B}_{1}{C}_{1}$ 中， ${AB} = A{A}_{1}$ ，点 $D$ 是${A}_{1}{B}_{1}$ 的中点,点 $E$ 在 ${A}_{1}{C}_{1}$ 上,且 ${DE} \bot  {AE}$ .
    \begin{enumerate}
        \item 证明：平面 ${ADE} \bot$ 平面 ${AC}{C}_{1}{A}_{1}$ ；
        \item 求直线 ${AD}$ 和平面 ${AB}{C}_{1}$ 所成角的大小.
    \end{enumerate}

\begin{figure}[h]
    \centering
    \includegraphics[width=0.25\linewidth]{bodyxb//src/iahoafuhwoushfaouwhe.png}
\end{figure}

    \begin{jieda}
    \begin{enumerate}

        \item 证明：由正三棱柱 $ {ABC} - {A}_{1}{B}_{1}{C}_{1} $ 的性质,可得 $ {AA}_{1}\bot $ 平面 $ {A}_{1}{B}_{1}{C}_{1} $ ,
因为 $ {DE} \subset $ 平面 $ {A}_{1}{B}_{1}{C}_{1} $ ,所以 $ A{A}_{1} \bot  {DE} $ ,
又因为 $ {DE}\bot {AE},{A{A}_{1}} \cap  {AE} = A $ 且 $ A{A}_{1},{AE} \subset $ 平面 $ {AC}{C}_{1}{A}_{1} $ ,
所以 $ {DE} \bot $ 平面 $ {AC}{C}_{1}{A}_{1} $ ,
因为 $ {DE} \subset $ 平面 $ {ADE} $ ,所以平面 $ {ADE} \bot $ 平面 $ {AC}{C}_{1}{A}_{1} $ .
        \item 线面角为$arcsin\dfrac{2\sqrt{105}}{35}$.

    \end{enumerate}
    
        \begin{figure}[h]
            \centering
            \includegraphics[width=0.25\linewidth]{bodyxb//src/aouehrgojahfrojahOjhfawfa.png}
        \end{figure}
    \end{jieda}
        


    \item 如图，四棱锥 $P - {ABCD}$ 中，底面 ${ABCD}$ 为矩形， ${PA} \bot$ 底面 ${ABCD},{PA} = {AB} = \sqrt{2}$ ，点 $E$ 是棱 ${PB}$ 的中点.
    \begin{enumerate}
        \item 证明： ${AE}\bot$ 平面 ${PBC}$ ；
        \item 若 ${AD} = 1$ ，求二面角 $B - {EC} - D$ 的大小.
    \end{enumerate}
    
    \begin{figure}[h]
        \centering
        \includegraphics[width=0.25\linewidth]{bodyxb//src/uiahf'auhrfoauhw'fuh.png}
    \end{figure}

    \begin{jieda}
        \begin{enumerate}
            \item 证明：由 $ {PA}\bot $ 平面 $ {ABCD} $ , $ {AB},{BC} \subset $ 平面 $ {ABCD} $ 可得 $ {PA}\bot {AB},{PA}\bot {BC} $ .又 $ {PA} = {AB} $ ,故 $ \Delta {PAB} $ 为等腰直角三角形,而 $ E $ 为棱 $ {PB} $ 的中点,所以 $ {AE}\bot {PB} $ ,由题意 $ {BC} \bot  {AB} $ ,而 $ {PA} \cap  {AB} = A $ ,故 $ {BC} \bot $ 平面 $ {PAB} $ ,而 $ {AE} \subset $ 平面 $ {PAB} $ ,故 $ {AE} \bot  {BC} $ ,而 $ {PB} \cap  {BC} = B $ ,故 $ {AE} \bot $ 平面 $ {PBC} $ ,因为 $ {AE} \subset $ 平面 $ {ADE} $ ,故平面 $ {ADE} \bot $ 平面 $ {PBC} $ .
            \item 二面角的大小为$\pi -arccos \dfrac{\sqrt{3}}{3}$ .
        \end{enumerate}
    \end{jieda}

    
    \item 如图,在 Rt $\bigtriangleup {ABC}$ 中, $\angle C = {90}^{ \circ  },{BC} = 3,{AC} = 6,D,E$ 分别是 ${AC},{AB}$ 上的点,且 ${DE} \; //{{BC}},{{DE}} = 2$ ,将 $\bigtriangleup {{ADE}}$ 沿 DE 折起到 $\bigtriangleup {{A}}_{1}{{DE}}$ 的位置,使 ${{A}}_{1}{C} \bot  {{CD}}$ ,如图 2.
    \begin{enumerate}
        \item 求证： ${{A}}_{1}{C}\bot$ 平面 ${{BCDE}}$ ；

        \item 若M是 ${{A}}_{1}{D}$ 的中点，求 ${{CM}}$ 与平面 ${{A}}_{1}{{DP}}$ 与平面 ${{A}}_{1}{{BE}}$ 垂直？说明理由.
        
        \item 线段BC上是否存在点P，使平面 ${{A}}_{1}{{DP}}$ 与平面 ${{A}}_{1}{{BE}}$ 垂直？说明理由.
    \end{enumerate}

\begin{figure}[h]
    \centering
    \includegraphics[width=0.5\linewidth]{bodyxb//src/isgjaoehrfaouehr.png}
\end{figure}

    \begin{jieda}
        \begin{enumerate}
            \item 证明：因为 $ {DE}//{BC},\angle C = {90}^{ \circ  } $ ,所以 $ {DE}\bot {AC} $ .所以 $ {CD} \bot  {DE},{A}_{1}D \bot  {DE} $ ,所以 $ {DE} \bot $ 平面 $ {A}_{1}{CD} $ .又因为 $ {A}_{1}C \subset $ 平面 $ {A}_{1}{CD} $ ,所以 $ {A}_{1}C \bot  {DE} $ ,又 $ {A}_{1}C \bot  {CD} $ ,所以 $ {A}_{1}C \bot $ 平面 $ {BCDE} $ .
            
            \item 解：如图,以 $ C $ 为坐标原点,建立如图所示空间直角坐标
            系 $ C - {xyz} $ ,则 $ D\left( {-2,0,0}\right) ,{A}_{1}\left( {0,0,2\sqrt{3}}\right) ,B\left( {0,3,0}\right) $ ,$ E\left( {-2,2,0}\right) $ . 所以 $ \vv{{A}_{1}B} = \left( {0,3, - 2\sqrt{3}}\right) ,\vv{{A}_{1}E} = \left( {-2,2, - 2\sqrt{3}}\right) $ . 设平面 $ {A}_{1}{BE} $ 法向量为 $ \vv{n} = \left( {x,y,z}\right) $ . 则 $ \left\{  \begin{array}{l} \vv{{A}_{1}B} \cdot  \vv{n} = 0, \\  \vv{{A}_{1}E} \cdot  \vv{n} = 0. \end{array}\right. $ 所以 $ \left\{  \begin{array}{l} {3y} - 2\sqrt{3}z = 0, \\   - {2x} + {2y} - 2\sqrt{3}z = 0. \end{array}\right. $ 则 则$ \left\{  \begin{array}{l} z = \frac{\sqrt{3}}{2}y, \\  x =  - \frac{y}{2}. \end{array}\right. $ 故 $ \vv{n} = \left( {-1,2,\sqrt{3}}\right) $ ,又因为 $ M\left( {-1,0,\sqrt{3}}\right) $ ,所以 $ \vv{CM} = \left( {-1,0,\sqrt{3}}\right) $ .设 $ {CM} $ 与平面 $ {A}_{1}{BE} $ 所成角的大小为 $ \theta $ ,则$ \sin \theta  = \left| {\cos \langle \vv{n},\vv{CM}\rangle }\right|  = \frac{\vv{CM} \cdot  \vv{n}}{\left| \vv{CM}\right| \left| \vv{n}\right| } = \frac{1 + 3}{\sqrt{1 + 3} \cdot  \sqrt{1 + 4 + 3}} = \frac{\sqrt{2}}{2}. $故 $ {CM} $ 与平面 $ {A}_{1}{BE} $ 所成角的大小为 $ {45}^{ \circ  } $ .

            \begin{figure}[h]
                \centering
                \includegraphics[width=0.5\linewidth]{bodyxb//src/fALhfaohwefouahwofa.png}
            \end{figure}

            \item 解：线段 $ {BC} $ 上不存在点 $ P $ ,使平面 $ {A}_{1}{DP} $ 与平面 $ {A}_{1}{BE} $ 垂直. 理由如下：假设这样的点 $ P $ 存在,其坐标为(0,a,0),其中 $ a \in  \left\lbrack  {0,3}\right\rbrack $ ,则 $ \vv{{A}_{1}P} = \left( {0,a, - 2\sqrt{3}}\right) ,\vv{DP} = \left( {2,a,0}\right) $ .设平面 $ {A}_{1}{DP} $ 法向量为 $ \vv{{n}_{1}} = \left( {{x}_{1},{y}_{1},{z}_{1}}\right) $ ,则 $ \left\{  \begin{array}{l} a{y}_{1} - 2\sqrt{3}{z}_{1} = 0, \\  2{x}_{1} + a{y}_{1} = 0, \end{array}\right. $ 解得 $ \left\{  \begin{array}{l} {z}_{1} = \frac{\sqrt{3}}{6}a{y}_{1}, \\  {x}_{1} =  - \frac{1}{2}a{y}_{1}. \end{array}\right. $令 $ {y}_{1} = 6 $ ,于是有 $ \vv{{n}_{1}} = \left( {-{3a},6,\sqrt{3}a}\right) $ . 而平面 $ {A}_{1}{DP} $ 与平面 $ {A}_{1}{BE} $ 垂直,当且仅当 $ \vv{{n}_{1}} \cdot  \vv{n} = 0 $ ,即$ {3a} + {12} + {3a} = 0 $ ,解得 $ a =  - 2 $ ,与 $ a \in  \left\lbrack  {0,3}\right\rbrack $ 矛盾,故线段 $ {BC} $ 上不存在点 $ P $ ,使平面 $ {A}_{1}{DP} $ 与平面$ {A}_{1}{BE} $ 垂直.

        \end{enumerate}
    \end{jieda}

\end{enumerate}





\subsection*{空间向量在立体几何中的应用(2)}

A 组：

\begin{enumerate}

    \item 已知空间中三点 $A\left( {2,0,0}\right)$ ， $B\left( {0,2,0}\right)$ ， $C\left( {2,2,2}\right)$ ，则点 $C$ 到直线 ${AB}$ 的距离为\blkda{$\sqrt{6}$}.

    \item 已知法向量为 $\vv{n} = \left( {1,0,1}\right)$ 的平面 $\alpha$ 内有一点 $A\left( {-1,1,0}\right)$ ，则平面外点 $P\left( {-1,1,1}\right)$ 到平面 $a$ 的距离为\blkda{$\dfrac{\sqrt{2}}{2}$}.

    \item 在正方体 ${ABCD} - {A}_{1}{B}_{1}{C}_{1}{D}_{1}$ 中， ${AB} = 4$ ，则异面直线 ${AB}$ 和 ${A}_{1}C$ 的距离为\blkda{$2\sqrt{2}$}.

    \item 四面体 ${ABCD}$ 的所有棱长均为 2，则二面角 $A - {BC} - D$ 的大小为\blkda{$arccos\dfrac{1}{3}$}.


\end{enumerate}

B 组:

\begin{enumerate}

\item 平面 $a$ 经过点 $B\left( {1,0,0}\right)$ ，且 $a$ 的法向量 $\vv{n} = \left( {1,1,1}\right)$ ，则 $P\left( {1,2,3}\right)$ 到平面 $a$ 的距离为\blkda{$\dfrac{5\sqrt{3}}{3}$}.

\item 如图，正方体 ${ABCD} - {A}_{1}{B}_{1}{C}_{1}{D}_{1}$ 的棱长为 1，则点 ${A}_{1}$ 到平面 $B{C}_{1}D$ 的距离为\blkda{$\dfrac{2\sqrt{2}}{3}$}.

\begin{figure}[h]
    \centering
    \includegraphics[width=0.25\linewidth]{bodyxb//src/foaihfha'wuhwfaouwhefrw.png}
\end{figure}


\item 在三棱锥 $P - {ABC}$ 中， ${PA} \bot$ 底面 ${ABC}$ ， $D$ 是 ${PC}$ 的中点，已知$\angle {BAC} = \frac{\pi }{2},{AB} = 2,{AC} = 2\sqrt{3},{PA} = 2$ ，则异面直线 ${BC}$ 与 ${AD}$ 所成角的余弦值为\blkda{$\dfrac{3}{4}$}.

\item 已知平面 $a$ 的一个法向量 $\vv{n} = \left( {1,\sqrt{3},2}\right)$ ，直线 $l$ 的方向向量 $\vv{v} = \left( {1,0, - 1}\right)$ ，则直线 $l$ 与平面 $a$ 所成角的正弦值为\blkda{$\dfrac{1}{4}$}.


\item 已知直线 $l$ 的一个方向向量为 $\vv{a} = \left( {1, - 1,1}\right)$ ，平面 $a$ 的一个法向量为 $\vv{b} = \left( {1,0, - 1}\right)$ ，则直线 $l$ 与平面 $a$ 所成的角的大小为\blkda{0}.


\item 如图，在正四面体 $A - {BCD}$ 中， $\vv{CE} = \frac{2}{3}\vv{CD}$ ，则异面直线 ${AE}$ 与 ${BD}$ 所成角的余弦值为\blkda{$\dfrac{\sqrt{7}}{7}$}.

\begin{center}
    \includegraphics[width=0.25\linewidth]{bodyxb//src/faouhoufhawoaefa.png}
\end{center}


\item 若正四面体 $A - {BCD}$ 的侧面 ${ABC}$ 内一动点 $P$ 到底面 ${BCD}$ 的距离与到棱 ${AB}$ 的距离相等， 则动点 $P$ 的轨迹与 ${ABC}$ 组成图形可能是\dotfill{(\kh{D})}
\begin{center}
    \includegraphics[width=1.1\linewidth]{bodyxb//src/foaehfoauwhefouawe.png}
\end{center}

\begin{jieda}[h]
\begin{figure}
    \centering
    \includegraphics[width=0.99\linewidth]{bodyxb//src/faoihfoauhwbofjawebfoa.png}
\end{figure}

\end{jieda}

\item 正方体 ${ABCD} - {{A}^{\prime }{B}^{\prime }{C}^{\prime }{D}^{\prime }}$ 中， $E$ 、 $F$ 分别是 ${A}^{\prime }{C}^{\prime }$ 和 $A{B}^{\prime }$ 的中点,求证: 直线 ${EF}$ 平行于平面 ${BC}{C}^{\prime }{B}^{\prime }$ .

\begin{center}
    \includegraphics[width=0.25\linewidth]{oauhfoahwefohawjehfa.png}
\end{center}

\begin{jieda}
    取$BB'$的中点N，$B'C'$的中点M,可以证明$EF//MN$，从而由线线平行推线面平行.
\end{jieda}

\item 已知向量 $\vv{n} = \left( {k,4, - 2}\right)$ 是平面 $a$ 的法向量，单位向量 $\vv{m} = \left( {\frac{2}{3},s,t}\right)$ 是平面 $\beta$ 的法向量,平面 $a$ 平行于平面 $\beta$ ,求实数 $S$ 的值.
\begin{jieda}
    由$\vv{m}//\vv{n}$，可求得$s=\pm\dfrac{2}{3}$
\end{jieda}


\item 如图，正四棱柱 ${ABCD} - {A}^{\prime }{B}^{\prime }{C}^{\prime }{D}^{\prime }$ 中， ${AB} = 2,A{A}^{\prime } = 3,E$ 是 ${B}^{\prime }{C}^{\prime }$ 的中点，问在棱 ${D{D}^{\prime }}$ 上是否存在一点 $F$ ,使得 ${C}^{\prime }F$ 与平面 ${AE}{D}^{\prime }$ 平行? 若存在,求 ${D}^{\prime }F$ 的长; 若不存在, 请说明理由.
\begin{center}
    \includegraphics[width=0.25\linewidth]{bodyxb//src/aijfoahwojfajowfa.png}
\end{center}
\begin{jieda}
    存在，$D'F=\dfrac{1}{2}DD'=\dfrac{3}{2}$.
\end{jieda}

\item 确定下列直线与平面的位置关系，如果相交，求它们所成角的大小：
\begin{enumerate}
    \item 直线的一个方向向量为 $\left( {1, - 2,9}\right)$ ，平面的一个法向量为 $\left( {3, - 4,7}\right)$ ；
    \item 直线的一个方向向量为 $\left( {\frac{\sqrt{2}}{2}, - 1,\sqrt{2}}\right)$ ，平面的一个法向量为 $\left( {-1,\sqrt{2}, - 2}\right)$ .
\end{enumerate}

\begin{jieda}
    \begin{enumerate}
        \item $sin\alpha=\dfrac{\sqrt{1591}}{43}$，因此线面所成角为$arcsin\dfrac{\sqrt{1591}}{43}$；
        \item 因为$\vv{n}//\vv{d}$，所以角为$\dfrac{\pi}{2}$.
    \end{enumerate}
\end{jieda}

\item 在长方体 ${{ABCD}} - {{A}}_{1}{{\;B}}_{1}{{C}}_{1}{{D}}_{1}$ 中， ${{AB}} = 4,{{AD}} = 6,{{{AA}}}_{1} = 4$ ， ${M}$ 是 ${{A}}_{1}{{C}}_{1}$ 的中点，P 在线段 BC 上，且 ${{ICPI}} = 2,{Q}$ 是 ${{{DD}}}_{1}$ 的中点,求: \begin{enumerate}
    \item 异面直线 ${AM}$ 与 ${PQ}$ 所成角的余弦值; \item $M$ 到直线 PQ 的距离；
    \item M 到平面 ${{{AB}}}_{1}{P}$ 的距离。
\end{enumerate}

\begin{center}
    \includegraphics[width=0.25\linewidth]{bodyxb//src/akushfahfhawefa.png}
\end{center}

\begin{jieda}
    \begin{enumerate}
        \item $arccos\dfrac{\sqrt{174}}{58}$;
        \item $\dfrac{\sqrt{462}}{6}$;
        \item $\dfrac{5\sqrt{3}}{3}$
    \end{enumerate}
\end{jieda}



\item 如图，已知正三棱柱 ${ABC} - {A}_{1}{B}_{1}{C}_{1}$ 的底面边长为 8,对角线 ${B}_{1}C \; = {10},D$ 是 ${AC}$ 的中点。
\begin{enumerate}
    \item 求点 ${B}_{1}$ 到直线 ${AC}$ 的距离；
    \item 求直线 $A{B}_{1}$ 到面 ${C}_{1}{BD}$ 的距离。
\end{enumerate}

\begin{center}
    \includegraphics[width=0.25\linewidth]{bodyxb//src/o'aihfouarfjawjebrfa.png}
\end{center}

\begin{jieda}
    \begin{enumerate}
        \item $2\sqrt{21}$;
        \item $\dfrac{12\sqrt{13}}{13}$
    \end{enumerate}
\end{jieda}



\item 如图，在直三棱柱 ${ABC} - {A}_{1}{B}_{1}{C}_{1}$ 中， ${A{A}_{1}} = {AB} = 2$ ， ${AC} = 1$ ， $\angle {ACB} = {90}^{ \circ  }$ ， $D$ 是棱 ${AB}$ 上的一点.
\begin{enumerate}

    \item 若 ${AD} = {DB}$ ，求异面直线 ${B}_{1}D$ 与 $A{C}_{1}$ 所成的角的大小；
    
    \item ${CD} \bot  {B}_{1}D$ ，求点 $B$ 到平面 ${B}_{1}{CD}$ 的距离.
\end{enumerate}

\begin{center}
    \includegraphics[width=0.25\linewidth]{bodyxb//src/oahfoauhw'rifhbaerf.png}
\end{center}

\begin{jieda}
    \begin{enumerate}
        \item 由题设 $ C{C}_{1} \bot $ 面 $ {ABC} $ ,且 $ {AC}\bot {BC} $ , $ {A{A}_{1}} = {AB} = 2 $ , $ {AC} = 1 $ ,
所以 $ {BC} = \sqrt{4 - 1} = \sqrt{3} $ ,而 $ {AD} = {DB} = 1 $ ,
建如下图示空间直角坐标系,
则 $ C\left( {0,0,0}\right) ,A\left( {1,0,0}\right) ,B\left( {0,\sqrt{3},0}\right) ,{B}_{1}\left( {0,\sqrt{3},2}\right) ,{C}_{1}\left( {0,0,2}\right) ,D{\left( \frac{1}{2},\frac{\sqrt{3}}{2},0\right) }^{\prime } $
所以 $ \vv{{B}_{1}D} = \left( {\frac{1}{2}, - \frac{\sqrt{3}}{2}, - 2}\right) ,\vv{A{C}_{1}} = \left( {-1,0,2}\right) $ ,
则 $ \left| {\cos \left\langle  {\vv{{B}_{1}D},\vv{A{C}_{1}}}\right\rangle  }\right|  = \left| \frac{\vv{{B}_{1}D} \cdot  \vv{A{C}_{1}}}{\left| \vv{{B}_{1}D}\right| \left| \vv{A{C}_{1}}\right| }\right|  = \frac{\frac{9}{2}}{\sqrt{5} \times  \sqrt{5}} = \frac{9}{10} $ ,
所以异面直线 $ {B}_{1}D $ 与 $ A{C}_{1} $ 所成的角的大小为 $ \arccos \frac{9}{10} $ .
\begin{figure}[h]
    \centering
    \includegraphics[width=0.25\linewidth]{bodyxb//src/fa'ihfajbfwqoehfbohabe.png}
\end{figure}

        \item 若 $ \vv{BD} = \lambda \vv{BA} $ ,且 $ 0 \leq  \lambda  \leq  1 $ ,则 $ \vv{{B}_{1}D} = \vv{{B}_{1}B} + \vv{BD} = \vv{{B}_{1}B} + \lambda \vv{BA} = \left( {\lambda , - \sqrt{3}\lambda , - 2}\right) $ ,
所以 $ D = \left( {\lambda ,\sqrt{3} - \sqrt{3}\lambda ,0}\right) $ ,故 $ \vv{CD} = \left( {\lambda ,\sqrt{3} - \sqrt{3}\lambda ,0}\right) $ ,
由 $ {CD} \bot  {B}_{1}D $ ,则 $ \vv{CD} \cdot  \vv{{B}_{1}D} = {\lambda }^{2} - {3\lambda }\left( {1 - \lambda }\right)  = 0 $ ,可得 $ 4{\lambda }^{2} - {3\lambda } = 0 $ ,所以 $ \lambda  = 0 $ 或 $ \lambda  = \frac{3}{4} $ ,
当 $ \lambda  = 0 $ ,则 $ D = \left( {0,\sqrt{3},0}\right) $ 与 $ B $ 点重合,此时点 $ B $ 到平面 $ {B}_{1}{CD} $ 的距离为 0 ；
当 $ \lambda  = \frac{3}{4} $ ,则 $ D = \left( {\frac{3}{4},\frac{\sqrt{3}}{4},0}\right) $ ,则 $ \vv{DB} = \left( {-\frac{3}{4},\frac{3\sqrt{3}}{4},0}\right) $ , $ \vv{CD} = \left( {\frac{3}{4},\frac{\sqrt{3}}{4},0}\right) $ , $ \vv{C{B}_{1}} = \left( {0,\sqrt{3},2}\right) $ ,
令面 $ {B}_{1}{CD} $ 的一个法向量为 $ \vv{m} = \left( {x,y,z}\right) $ ,则 $ \left\{  \begin{array}{l} \vv{m} \cdot  \vv{CD} = \frac{3}{4}x + \frac{\sqrt{3}}{4}y = 0 \\  \vv{m} \cdot  \vv{C{B}_{1}} = \sqrt{3}y + {2z} = 0 \end{array}\right. $ ,
若 $ y =  - \sqrt{3} $ ,则 $ \vv{m} = \left( {1, - \sqrt{3},\frac{3}{2}}\right) $ ,
此时点 $ B $ 到平面 $ {B}_{1}{CD} $ 的距离为 $ \frac{\left| \vv{m} \cdot  \vv{DB}\right| }{\left| \vv{m}\right| } = \frac{3}{\sqrt{1 + 3 + \frac{9}{4}}} = \frac{6}{5} $ .
综上,点 $ B $ 到平面 $ {B}_{1}{CD} $ 的距离为 0 或 $ \frac{6}{5} $ .
    \end{enumerate}
\end{jieda}

\end{enumerate}


C组：
\begin{enumerate}
    \item 已知正四棱柱 ${{A}}_{1}{{A}}_{2}{{A}}_{3}{{A}}_{4} - {{B}}_{1}{{B}}_{2}{{B}}_{3}{{B}}_{4}$ 的底面为边长为 2，高为 3，则集合 $\left\{  {{x} \mid  {x} = {{A}}_{1}{{B}}_{1} \cdot  {{A}}_{i}{{B}}_{j},{i} \in  \{ 1,2,3,4\} ,{j} \in  \{ 1,2,3,4\} }\right\}$ 中元素的个数为\blkda{1}个。


    \item 定义一个集合 ${\Omega }$ ,集合中的元素是空间内的点集,任取 ${{P}}_{1},{{P}}_{2},{{P}}_{3} \in  \Omega$ , 存在不全为 0 的实数 ${\lambda }_{1},{\lambda }_{2},{\lambda }_{3}$ ,使得 ${\lambda }_{1}{\vv{OP}}_{1} + {\lambda }_{2}{\vv{OP}}_{2} + {\lambda }_{3}{\vv{OP}}_{3} = \vv{0}$ . 已知 $\left( {{1},{0},{0}}\right)  \in  \Omega$ ，则 $\left( {{0},{0},{1}}\right)  \notin  \Omega$ 的充分条件是\dotfill{(\kh{C})}

    \xx{$\left( {{0},{0},{0}}\right)  \in  {\Omega }$ } 
    {$\left( {-{1},{0},{0}}\right)  \in  {\Omega }$}
    {$\left( {{0},{1},{0}}\right)  \in  {\Omega }$}   
    {$\left( {{0},{0}, - {1}}\right)  \in  {\Omega }$}
    
\end{enumerate}
